% =====================================================================
%  CAPÍTULO 3 — REGISTRO DE SOFTWARE E EVOLUÇÃO DO PROJETO
% =====================================================================

\chapter{REGISTRO DE SOFTWARE E EVOLUÇÃO DO PROJETO}
\label{cap:registro-evolucao}

Este capítulo contextualiza o registro do \argus{} junto ao Instituto Nacional
da Propriedade Industrial (INPI) e documenta a evolução do projeto desde a
versão registrada até a versão atual apresentada neste trabalho. A apresentação
deixa explícita a \textbf{continuidade} entre as versões: o registro formalizou
a \textbf{versão inicial} do software, e o presente trabalho documenta o
aprimoramento e a maturação do sistema a partir dessa base.

\section{O registro de software no INPI}
\label{sec:registro-inpi}

O registro de programa de computador é um direito de propriedade intelectual
garantido pela Lei nº 9.609/1998, que dispõe sobre a proteção da propriedade
intelectual de programa de computador \cite{lei9609}. O registro é realizado no
INPI e assegura ao titular e aos autores a proteção jurídica sobre o código
fonte do software por um período de 50 anos a partir de 1º de janeiro do ano
subsequente à sua publicação \cite{inpi2026}.

O \argus{} foi registrado junto ao INPI, e o certificado oficial (exibido na
página de registro, no início deste documento) apresenta os dados
consolidados na Tabela~\ref{tab:registro}:

\begin{table}[h]
    \centering
    \caption{Dados do registro do software no INPI}
    \label{tab:registro}
    \small
    \begin{tabular}{>{\raggedright\arraybackslash}p{4.5cm}p{9cm}}
        \toprule
        \textbf{Campo} & \textbf{Valor} \\
        \midrule
        Número do processo & BR512026003786-7 \\
        Título & Argus Cyclist \\
        Data de criação / publicação & 20/01/2026 \\
        Titular & Fundação Universidade do Estado do Rio Grande do Norte (FUERN) \\
        Autores & Isaac de Lima Oliveira Filho; Paulo Sérgio Silva de Medeiros \\
        Linguagem & HTML; JavaScript; Outros \\
        Campo de aplicação & IF-10 \\
        Tipo de programa & SM-02 \\
        Algoritmo hash & SHA-512 \\
        Data de expedição & 02/06/2026 \\
        \bottomrule
    \end{tabular}
\end{table}

\section{A versão registrada como linha de base}
\label{sec:versao-inicial}

O registro do software foi realizado em 20/01/2026, momento em que o projeto se
encontrava em sua etapa inicial de desenvolvimento. A versão registrada ---
doravante denominada \textbf{versão inicial} --- contemplava as funcionalidades
essenciais que comprovaram a viabilidade técnica da proposta e definiram a
identidade do projeto:

\begin{itemize}
    \item \textbf{Simulador núcleo}: conexão BLE com sensores de frequência
        cardíaca e rolos de treino inteligentes (FTMS), renderização de mapa 3D
        e exportação de arquivos \texttt{.FIT}.
    \item \textbf{Histórico e gráficos}: banco de dados local (SQLite) com
        histórico de atividades e visualização gráfica de sessões.
    \item \textbf{Curva de potência (MMP)}: cálculo e exibição dos recordes
        de potência máxima em diferentes janelas de tempo.
    \item \textbf{Editor de rotas}: criação de rotas personalizadas
        diretamente no mapa.
\end{itemize}

Essa versão constituiu o \textit{baseline} funcional e arquitetural que foi
aprimorado ao longo do ciclo de vida do projeto.

\section{A evolução do \argus{} após o registro}
\label{sec:evolucao}

A partir da versão registrada, o projeto passou por um processo contínuo de
evolução orientado por \textit{sprints} de desenvolvimento incremental. Essa
evolução manifestou-se em três dimensões principais: (i) arquitetural e de
plataforma, (ii) funcional e (iii) de engenharia de qualidade.

\subsection{Evolução arquitetural e de plataforma}
\label{sec:evolucao-arquitetura}

A arquitetura da versão inicial, centrada no \textit{desktop}, foi ampliada
para suportar novos contextos de execução e novos padrões de projeto:

\begin{itemize}
    \item \textbf{Porta Android via Capacitor.js}: o \textit{frontend} passou a
        ser empacotado como aplicativo nativo Android, com um motor de
        simulação física reescrito em JavaScript para operação sem o
        \textit{backend} Go, permitindo que um smartphone comum atue como
        unidade de treino com conexão BLE nativa.
    \item \textbf{Migração do motor de mapa para MapLibre GL JS}: a renderização
        3D foi portada para uma solução de código aberto, com terreno
        topográfico (Terrain-RGB), projeção globo e temas dinâmicos.
    \item \textbf{Isolamento de dados por perfil}: a persistência evoluiu para
        um sistema multiperfil, com isolamento completo de dados entre atletas
        no mesmo dispositivo.
    \item \textbf{Padronização dos padrões de projeto}: consolidação dos
        padrões \textit{Repository} e \textit{Facade} na camada de persistência,
        com repositórios granulares por entidade.
    \item \textbf{Integrações externas}: além da exportação local de arquivos
        \texttt{.FIT}, foram adicionadas integrações OAuth 2.0 com o Strava e,
        posteriormente, com o Intervals.icu, com fila de sincronização
        assíncrona resiliente a falhas de rede.
\end{itemize}

\subsection{Evolução funcional}
\label{sec:evolucao-funcional}

A evolução funcional foi substancial e pode ser organizada em grandes blocos
cronológicos:

\begin{itemize}
    \item \textbf{Primeiro trimestre (fev.--abr. 2026)}: implementação de
        treinos estruturados com importação de planos ZWO, controle de
        intensidade, testes de aptidão (FTP e \textit{Ramp Test}), migração do
        mapa para MapLibre, integração com o Strava (autenticação OAuth 2.0 e
        \textit{upload} automático), modos de evento (Sprint, KOM e
        Contrarrelógio) com \textit{leaderboard} persistente e análises
        avançadas (PMC, TRIMP, HRR, desvio cardiovascular).
    \item \textbf{Segundo trimestre (mai.--jun. 2026)}: porta Android via
        Capacitor, modo Estúdio (\textit{Studio Mode}), sistema de
        gamificação expandido (troféus, conquistas de consistência e
        recuperação, metas personalizadas), assistente de treino com IA
        generativa local (Ollama) com geração de planos ZWO, e rastreamento de
        desgaste de componentes mecânicos da bicicleta.
    \item \textbf{Terceiro trimestre (jul.--ago. 2026)}: unificação dos modos de
        estúdio e desafio em um HUD compartilhado, modo \textit{arena} com
        classificação filtrada por gênero, integração com o Intervals.icu,
        correções de estabilidade e melhorias de confiabilidade na camada BLE.
\end{itemize}

\section{Comparativo entre a versão registrada e a versão atual}
\label{sec:comparativo}

A Tabela~\ref{tab:evolucao} resume as principais diferenças entre a versão
inicial registrada e a versão atual do software:

\begin{table}[h]
    \centering
    \caption{Evolução do \argus{}: versão registrada \textit{versus} versão atual}
    \label{tab:evolucao}
    \small
    \begin{tabular}{
        >{\raggedright\arraybackslash}p{3.6cm}
        >{\raggedright\arraybackslash}p{5.4cm}
        >{\raggedright\arraybackslash}p{5.4cm}
    }
        \toprule
        \textbf{Aspecto} & \textbf{Versão registrada (01/2026)} &
        \textbf{Versão atual (2026)} \\
        \midrule
        Plataforma & Desktop (Wails) & Desktop + Android (Capacitor) \\
        Motor de mapa & Mapbox 3D & MapLibre GL JS (Terrain-RGB) \\
        Motor físico & Go (backend) & Go + portado para JavaScript \\
        Conectividade BLE & FTMS + HRM & FTMS, FE-C, HRM, potência, \textit{virtual trainer} \\
        Treinos & Livres & Livres + estruturados (ZWO) + testes FTP \\
        Métricas & Básicas & NP, IF, TSS, TRIMP, PMC, HRR, desvio cardiovascular \\
        Integrações & --- & Strava (OAuth 2.0) e Intervals.icu \\
        IA & --- & Assistente local via Ollama (planos ZWO) \\
        Gamificação & --- & XP, níveis, troféus, conquistas, metas \\
        Modos de treino & --- & Estúdio, Sprint, KOM, Contrarrelógio \\
        Manutenção & --- & Rastreamento de desgaste de componentes \\
        Persistência & Perfil único & Multiperfil isolado (SQLite) \\
        \bottomrule
    \end{tabular}
\end{table}

\section{Continuidade entre as versões}
\label{sec:continuidade}

É importante destacar que a evolução descrita não representa o abandono da
versão registrada, mas sim a sua \textbf{maturação}. As funcionalidades
essenciais da versão inicial --- simulação BLE, mapa 3D, exportação
\texttt{.FIT} e histórico local --- foram mantidas, aprimoradas e expandidas, e
a arquitetura original orientou todas as decisões posteriores de projeto. O
registro no INPI, portanto, protege a base do software, enquanto este documento
formaliza o estado atual do sistema, evidenciando o caráter contínuo e
incremental do desenvolvimento do \argus{}.
